The calibration loop described in \hyperref[subsec:calibration]{Calibration} operates 
at the level of the full interaction. What follows describes the finer 
movement within it --- the emotional sequence that tends to unfold naturally 
when presence is genuine and neither person is managing toward an outcome.

This is not a protocol to execute. It is a description of what the atmosphere 
looks like when the loop is working honestly.

\vspace{0.5em}
\textbf{Register Matching.}
Before anything else: feel what is actually in the room. Not what you want 
to produce --- what is already there. Her current emotional state, her 
register, the quality of her attention. The first movement is not transmission. 
It is reception. Emotional contagion operates in both directions; the shared 
field begins forming the moment one person genuinely attends to the other. 
Baseline First (\hyperref[subsec:baseline]{Baseline}) is the prerequisite. 
You cannot match a register you have not read.

\vspace{0.5em}
\textbf{Synchrony.}
When attention is genuinely outward and neither person is performing, 
the emotional states of two people begin to converge without either directing 
it. This is not mirroring --- mirroring is a technique. Synchrony is a 
consequence. The shared field has formed when her expression begins arriving 
from the same space yours is arriving from: incomplete thoughts, 
unhurried silences, the same quality of attention moving in both directions. 
Neither person caused this. Both people allowed it.

\vspace{0.5em}
\textbf{Gradual Introduction.}
Once synchrony is present, what is real in the speaker becomes naturally 
less contained. Not because he decided to reveal more --- because the shared 
field makes containment less necessary. This is where Graze and Gradual 
(\hyperref[subsec:three_modes]{three modes}) enter: something genuine is noticed, 
lightly expressed, and left where it lands. Not pursued. Not explained. 
Just placed --- and then moved past.

What this produces in her is not a response to a move. It is a response 
to an atmosphere. She begins contributing her own fragments, her own 
hesitations, her own unguarded moments --- not because she was invited to, 
but because the space made it natural. This is the movement into 
co-authorship. Not a transition that was engineered. A threshold that 
was crossed because neither person was guarding it.

\vspace{0.5em}
The sequence is: \textit{receive before transmitting. Allow synchrony 
rather than producing it. Let what is real become less contained 
as the shared field deepens. Do not cross into co-authorship. 
Arrive there.}

\begin{calloutbox}
Feel her register before you transmit yours.\\
Enter the same space before you introduce your own.\\
Introduce your own lightly, and move past it.\\[0.4em]
What follows is not something you produce.\\
It is something you allow.
\end{calloutbox}

\vspace{0.5em}
\textbf{Register Congruence.}
Matching her register is not accommodation. It is accuracy.
If the shared field is dry, warm expression does not balance it — it breaks it.
What the room produced is the honest material.
Expressing from a different temperature imports something the field did not contain.

This is not mimicry. Mimicry is downstream and deliberate —
a decision to reflect what was observed.
Register congruence is upstream: the shared field has already shaped
what forms before the choice to speak.
The man who stays warm when the room is dry
is not being himself. He is failing to notice
that the room has already moved through him —
and that what it left behind is the honest starting point.

The push and pull that follows is not exchange of energy.
It is rhythm — back and forth in pacing and response,
not in emotional temperature.
The temperature belongs to the field.
The rhythm belongs to the two people inside it.
One is condition. The other is movement within it.